% Paper 4, §6 — Truth directions and the geometry of truth
% Thesis: truth is the most-studied probing target since 2022. The field
% has converged on the claim that truth is linearly decodable from
% late-mid layers of LLaMA-family models on factual statements; other
% claims (cross-dataset alignment, CCS-under-negation, single-axis vs.
% multi-axis) remain open. Marks-Tegmark 2023 "stark misalignment" and
% Vennemeyer 2025 sycophancy decomposition generalize a single pattern:
% behaviorally-singular properties may be representationally-plural.

\section{Truth directions and the geometry of truth}
\label{sec:truth-directions}

If \S\ref{sec:probing} sets up the probing methodology and
\S\ref{sec:correlation-causation} closes the correlation--causation gap
with activation interventions, this section is the worked case study
that exercises both. ``Truth'' --- whether the model's \glsterm{residual-stream}{residual stream}
encodes the truth \glsterm{value}{value} of a factual statement as a linearly decodable
direction --- is the most-studied probing target since 2022, and the
empirical record is rich enough to separate what the field has actually
converged on from what remains open. The convergence claim is narrow:
on the LLaMA-family on simple factual true/false statements, a linear
direction in the \glsterm{residual-stream}{residual stream} at late-mid layers carries the truth
\glsterm{value}{value} with high \glsterm{probe-2}{probe} accuracy and is causally implicated in the
model's TRUE/FALSE prediction \citep{marks-tegmark-2023-truth}. The
open claims are the load-bearing ones: whether the \emph{same}
direction is recovered across prompt formats, whether the
unsupervised CCS direction recovers truth or a near-truth confound,
whether ``truth'' is a single linear axis or a multi-axis subspace,
and whether any of this transfers cross-model-family. We work through
each in turn.

\subsection{The Marks--Tegmark convergence claim}
\label{sec:truth-marks-tegmark}

\citet{marks-tegmark-2023-truth} make the strongest published case for
a linear truth representation, with three lines of evidence assembled
on LLaMA-2 7B/13B/70B and the curated \texttt{cities} /
\texttt{neg\_cities} / \texttt{larger\_than} / \texttt{cities\_conj}
datasets of \S\ref{sec:probe-protocol}\footnote{KB excerpt:
\texttt{kb/excerpts/marks-tegmark-2023-truth.md\#abstract}.}:
\emph{visualization} (PCA of layer-$\ell$ activations shows clear
linear true/false separation in the top two PCs); \emph{transfer}
(probes trained on one dataset classify a structurally distinct
dataset above chance and ablation-flip its predictions); and
\emph{causal intervention} (ablating the \glsterm{probe-2}{probe}-derived direction in
the \glsterm{residual-stream}{residual stream} during the forward pass causes the model to
``treat false statements as true and vice versa''
\citep[abstract]{marks-tegmark-2023-truth}). The headline
methodological recommendation is that, among the four \glsterm{probe-2}{probe} families
of \S\ref{sec:probing}, the parameter-free \emph{mass-mean} direction
$\mathbf{d}_\mathrm{MM} = \boldsymbol{\mu}_+ - \boldsymbol{\mu}_-$ of
Eq.~\eqref{eq:mass-mean} is the one ``most causally implicated in
model outputs''~\citep[\S 1]{marks-tegmark-2023-truth} --- generalizes
across datasets at least as well as logistic regression, while
intervening on it produces larger behavioral effects.

The localization picture is sharp. Causal patching localizes truth
representations to a small group of hidden states ``over the final
token of the statement and end-of-sentence punctuation,'' which the
authors hypothesize ``store a representation of the statement's
truth''~\citep[\S 3]{marks-tegmark-2023-truth}\footnote{KB excerpt:
\texttt{kb/excerpts/marks-tegmark-2023-truth.md\#sec-3-layers}.}. On
LLaMA-2-13B the separation emerges in early-middle layers (around
$\ell = 12$), is salient through the rest of the stack, and ``emerges
later for datasets of more structurally complex statements (e.g.\
conjunctive statements)''~\citep[\S 4]{marks-tegmark-2023-truth}. The
pattern --- linear separation, localized to a layer band, summarized
at the period-token --- is the picture the field now defaults to when
saying ``the \glsterm{llm}{LLM} has a \glsterm{truth-direction-2}{truth direction}.''

\begin{intuition}
Mass-mean turns a probing claim into a claim about the data's class
geometry: $\mathbf{d}_\mathrm{MM} = \boldsymbol{\mu}_+ -
\boldsymbol{\mu}_-$ is the axis along which the per-class means are
separated, and the projection $\mathbf{h}^\top \mathbf{d}_\mathrm{MM}$
is the ``how far does this activation lean toward the true class''
scalar. Why does the parameter-free direction beat trained logistic
regression on \emph{causal} mediation? Because logistic regression
fits the discriminative direction --- the axis that minimizes
classification error --- and that axis can read a confound $z$ that
linearly separates the data without the model using $z$. Mass-mean
fits no boundary; it points along the class-difference of the data
itself, so the recovered axis depends only on $\mathcal{D}$, not on
an optimizer. Returning to canonical math: with isotropic Gaussian
classes of equal covariance, mass-mean and logistic regression
recover the \emph{same} direction up to scale (Fisher LDA); under
heteroscedastic or non-Gaussian residuals they diverge, and the
divergence is exactly the discriminative flexibility that mass-mean
chooses to forgo.
\end{intuition}

\subsection{The Burns CCS direction --- truth-like, but unsupervised}
\label{sec:ccs-truth}

Contrast-Consistent Search (CCS)
\deepencite{burns2022-ccs §3, citation-pending} attacks the same
question from the unsupervised side. Given an unlabeled set of
contrast pairs $(\mathbf{x}^+, \mathbf{x}^-)$ --- e.g., a statement
appended with ``yes'' versus ``no'' --- CCS fits a \glsterm{linear-probe}{linear probe}
$p_\theta(\mathbf{h}) = \sigma(\mathbf{w}^\top \mathbf{h} + b)$ by
minimizing a \glsterm{self-consistency-2}{self-consistency} loss
\begin{equation}
  \mathcal{L}_\mathrm{CCS}(\theta)
  \;=\;
  \mathbb{E}_{(\mathbf{x}^+, \mathbf{x}^-)}\!\left[
    \big(p_\theta(\mathbf{h}^+) + p_\theta(\mathbf{h}^-) - 1\big)^2
    \;+\;
    \min\!\big(p_\theta(\mathbf{h}^+),\, p_\theta(\mathbf{h}^-)\big)^2
  \right],
  \label{eq:ccs-loss}
\end{equation}
which jointly enforces (i) that the \glsterm{probe-2}{probe} assigns complementary
probabilities to the two halves of the pair --- a constraint any
\emph{truth} \glsterm{feature}{feature} must satisfy, since a statement and its negation
have opposite truth values --- and (ii) that the \glsterm{probe-2}{probe} is confident
on at least one side. The discovered $\mathbf{w}$ is reported to be
``truth-like'': \glsterm{probe-2}{probe} accuracy on held-out labeled true/false data is
competitive with supervised logistic regression on in-distribution
splits \deepencite{burns2022-ccs §4, citation-pending}. The
contribution is conceptual: CCS removes the curated-label confound
that motivates Step 1 of the probing protocol of
\S\ref{sec:probe-protocol}; the discovered direction is a property of
the model's representation \emph{and} of the contrast-pair templates,
not of any human-curated label.

The empirical record on CCS, however, is more cautious than the
in-distribution accuracy headline suggests.
\citet{marks-tegmark-2023-truth} situate CCS alongside other probing
techniques in their related work and note that \emph{transfer to
negated statements} is the load-bearing
test\footnote{KB excerpt:
\texttt{kb/excerpts/marks-tegmark-2023-truth.md\#sec-1-related}.}: a
direction that separates ``X is true'' from ``X is false'' but fails
on ``X is not false'' versus ``X is not true'' is reading something
correlated with truth (sentiment, statement plausibility, surface
template features) rather than truth itself. The pattern in their
related-work discussion is that prior unsupervised and contrastive
probing work has documented ``failures of these probes to generalize
in basic ways''~\citep[abstract]{marks-tegmark-2023-truth} --- and the
\texttt{cities} versus \texttt{neg\_cities} contrast is exactly the
form of generalization test on which the
\citeauthor{marks-tegmark-2023-truth} mass-mean direction does pass
and on which CCS-style directions are documented to struggle. We
treat this as a strong empirical signal that the unsupervised
contrast-pair objective of Eq.~\eqref{eq:ccs-loss} can latch onto a
near-truth confound rather than truth proper, and we flag the
underlying claim as anchored to in-distribution accuracy until a
peer-reviewed cross-dataset reproduction is in hand.

\subsection{Stark misalignment: the dataset-dependent direction}
\label{sec:truth-misalignment}

The cleanest open question is the one
\citet{marks-tegmark-2023-truth} surface in their own \S 4 visualization
analysis. After establishing that each of \texttt{cities},
\texttt{neg\_cities}, \texttt{larger\_than}, and \texttt{cities\_conj}
admits a linearly separating direction, the authors ask whether those
directions \emph{align} across datasets. The answer is: they align
\emph{at late layers}, but not at early/mid
layers\footnote{KB excerpt:
\texttt{kb/excerpts/marks-tegmark-2023-truth.md\#sec-4-misalignment}.}.
The verbatim characterization is worth quoting in full:

\begin{quote}
\itshape
The axes of separation for various true/false datasets align often,
but not always. \dots\ Fig.~3(c) shows stark failures of alignment,
with the axes of separation for datasets of statements and their
negations being approximately
orthogonal.~\citep[\S 4]{marks-tegmark-2023-truth}
\end{quote}

\noindent and the layer-wise dynamics:

\begin{quote}
\itshape
In early layers, \texttt{cities} and \texttt{neg\_cities} separate
antipodally, before rotating to lie orthogonally\dots, and finally
aligning in later layers.~\citep[\S 4 / App.~C]{marks-tegmark-2023-truth}
\end{quote}

The implication for the ``does the model have a \glsterm{truth-direction-2}{truth direction}''
question is direct. The \glsterm{probe-2}{probe} direction recovered from
\texttt{cities} alone is \emph{not} the \glsterm{probe-2}{probe} direction recovered from
\texttt{neg\_cities}, in early or middle layers; it would be a mistake
to treat either single-dataset-derived $\hat{\mathbf{d}}_y$ as a
canonical truth axis at those depths. Only at late layers do the
per-dataset directions begin to share a common subspace, and the
authors flag a scale dimension on top of the layer dimension:
\texttt{larger\_than} and \texttt{smaller\_than} separate
\emph{antipodally} in LLaMA-2-13B but along a \emph{common} direction
in LLaMA-2-70B~\citep[\S 4]{marks-tegmark-2023-truth}, suggesting that
larger models compress what 13B encodes as multiple per-template
axes into a single more-abstract axis.

\subsection{Vennemeyer 2025: the multi-component generalization}
\label{sec:vennemeyer-multicomponent}

\citet{vennemeyer2025-sycophancy-not-one-thing}
\deepencite{vennemeyer2025-sycophancy-not-one-thing §3, citation-pending}
extend the multi-component finding from truth to behavioral
phenomena. Applying difference-in-means probing to \glsterm{sycophancy}{sycophancy} on
LLaMA-family models, they report that the behavior decomposes into
\emph{three independently-steerable directions}: agreement (do I
match the user's stated position?), praise (do I compliment the
user?), and genuine-agreement (do I give the answer I would have
given without the user's lead?). Activation-steering along each
direction independently changes a different facet of the sycophantic
behavior; combining them at chosen scales reproduces the full
behavioral profile. The methodological generalization is direct:
where \citeauthor{marks-tegmark-2023-truth} document
\emph{representational plurality} for ``truth'' \emph{across datasets
within a model}, \citeauthor{vennemeyer2025-sycophancy-not-one-thing}
document representational plurality for \emph{a single behavior at a
single distribution} --- both observations point at the same
underlying claim: \textbf{behaviorally-singular properties may be
representationally-plural}, and a single direction is a
\emph{projection} of a richer subspace rather than the full
representation.

\begin{contradiction}
\textbf{Is there a canonical \glsterm{truth-direction-2}{truth direction}?} Three observations
strain the strong-form claim that ``the \glsterm{llm}{LLM} linearly represents
truth'' as a single axis $\mathbf{d}_\mathrm{truth} \in \mathbb{R}^D$.
\textsf{(1)} Within a model,
\citet[\S 4]{marks-tegmark-2023-truth} document early-/mid-layer
``stark misalignment'' --- per-dataset truth directions are
antipodal, rotate to orthogonal, then align only late: a
single-dataset-derived $\hat{\mathbf{d}}_y$ is not a depth-uniform
truth axis. \textsf{(2)} CCS, the unsupervised analog of
Eq.~\eqref{eq:ccs-loss}, is documented to suffer generalization
failures under negation \citep[\S 1.1]{marks-tegmark-2023-truth}
\deepencite{burns2022-ccs §5, citation-pending} --- a putative truth
direction that fails the most basic logical operation on truth is
better understood as a near-truth confound than as truth.
\textsf{(3)} The
\citet{vennemeyer2025-sycophancy-not-one-thing} \glsterm{sycophancy}{sycophancy}
decomposition shows that an alignment-relevant phenomenon thought to
be representationally-singular admits a multi-axis decomposition; the
generalization to truth is direct (and tested informally in their
follow-up). \textbf{What evidence would resolve it.} A multi-axis
basis recovery procedure that, applied independently to
\texttt{cities}, \texttt{neg\_cities}, \texttt{larger\_than}, and
\texttt{cities\_conj} at matched layers, returns a low-dimensional
subspace whose per-dataset projections recover the
single-dataset directions modulo a known equivalence (rotation,
sign, scale). Until then, ``the \glsterm{truth-direction-2}{truth direction}'' is a useful
abstraction at late layers in fixed model + dataset family --- a
\emph{representational subspace} treated as a one-dimensional
projection --- not a basis-free property of the model.
\end{contradiction}

\subsection{What the field has converged on, and what remains open}
\label{sec:truth-frontier}

The convergence picture, restricted to load-bearing claims that
multiple labs and probing methodologies have triangulated:

\begin{itemize}
  \item \textbf{Late-mid layer linear decodability.} On the LLaMA-2
    family, simple factual true/false statements admit a \glsterm{linear-probe}{linear probe}
    with high held-out accuracy in residual-stream activations at
    layers $\ell \approx 12\!-\!15$ (LLaMA-2-13B) and equivalent
    relative depths on 7B/70B; the direction remains salient through
    the rest of the stack \citep[\S 3, \S 4]{marks-tegmark-2023-truth}.
  \item \textbf{Period-token summarization.} The most-informative \glsterm{probe-2}{probe}
    position is the end-of-statement punctuation token; this
    \glsterm{summarization-behavior}{summarization behavior} was observed prior on sentiment
    \citep[\S 3, citing Tigges et al.\ 2023]{marks-tegmark-2023-truth}
    and the truth result reproduces the geometric pattern.
  \item \textbf{Mass-mean as the direction with strongest causal
    mediation.} Among the \glsterm{probe-2}{probe} families of \S\ref{sec:probing},
    $\mathbf{d}_\mathrm{MM}$ flips TRUE/FALSE under direction-ablation
    more strongly than logistic-regression directions on matched
    datasets \citep[\S 1, \S abstract]{marks-tegmark-2023-truth}.
  \item \textbf{Larger models' truth axes are more abstract.}
    Per-template direction differences shrink as model size grows;
    LLaMA-2-70B compresses what 13B represents as multiple per-template
    axes into a more shared axis~\citep[\S 4.1]{marks-tegmark-2023-truth}.
\end{itemize}

The open claims, in falsification-priority order:
\begin{itemize}
  \item \textbf{Cross-prompt-format alignment.} Whether the
    early-/mid-layer per-dataset misalignment is a deficiency of the
    \emph{\glsterm{probe-2}{probe}} (single-direction summary of a multi-axis subspace)
    or a property of the \emph{model} (no shared early-layer truth
    representation across prompt formats) is unresolved
    \citep[\S 4]{marks-tegmark-2023-truth}.
  \item \textbf{CCS under negation.} The empirical record on whether
    Eq.~\eqref{eq:ccs-loss} recovers truth versus a confound on
    negated-statement transfer is fragmented across follow-ups; a
    matched-protocol reproduction at LLaMA-2-13B/70B scale is missing
    \deepencite{burns2022-ccs §5, citation-pending}.
  \item \textbf{Multi-class generalization of mass-mean.}
    Eq.~\eqref{eq:mass-mean} is binary by construction; the analog
    for $K$-class truthfulness gradations (e.g., model-graded
    confidence buckets) is not yet a settled methodological recipe.
  \item \textbf{Cross-model-family transfer.} Truth directions in
    LLaMA-2-13B and Qwen-2.5-14B almost certainly do not have the
    same coordinates (the bases differ); whether there is a
    canonical \glsterm{truth-direction-2}{truth direction} up to model-specific basis or whether
    truth is fundamentally model-dependent is the question that
    Paper~5 needs answered to sharpen behavioral probes for
    deception, \glsterm{sycophancy}{sycophancy}, \glsterm{scheming}{scheming}, and alignment-faking into
    cross-model evaluation tools.
\end{itemize}

\subsection{Forward link}
\label{sec:truth-forward}

Behavioral probes for alignment-relevant phenomena ride on the
methodology this section assembles. \hyperref[P5-sec:sycophancy]{Paper~5 §9} (\glsterm{sycophancy}{sycophancy})
consumes the
\citeauthor{vennemeyer2025-sycophancy-not-one-thing} multi-component
result directly; \hyperref[P5-sec:scheming]{Paper~5 §10} (\glsterm{scheming}{scheming}) and \hyperref[P5-sec:alignment-faking]{§11} (alignment-faking) use
the same \glsterm{probe-2}{probe}-then-causally-validate pipeline of
\S\ref{sec:correlation-causation} with behavioral targets in place of
factual ones. The single methodological lesson the truth case study
exports is the contradiction box of \S\ref{sec:truth-misalignment}:
treat the recovered direction as a single-dataset summary, not as a
basis-free property; report cross-dataset transfer in
\emph{both} accuracy \emph{and} ablation effect; and prefer multi-axis
subspace recovery over single-axis fits when the underlying behavior
admits a multi-component decomposition.

The same canonicality question generalizes one floor down, to the
\emph{unsupervised} \glsterm{feature}{feature} framing. Sparse autoencoders
(\S\ref{sec:sparse-autoencoders}, written) decompose residual-stream
activations into a learned dictionary of $M \gg D$ directions, each
plausibly a ``\glsterm{feature}{feature}'' candidate; whether two \glsterm{sparse-autoencoder}{SAE} training runs at
different sparsity targets recover the \emph{same} \glsterm{feature}{feature} inventory
is the structural cousin of ``do two probes on different truth datasets
recover the same direction.''
\S\ref{sec:sae-canonicality} takes that question up at the \glsterm{sparse-autoencoder}{SAE} level,
and the cross-method synthesis of \S\ref{sec:cross-method} returns to
both: \glsterm{probe-2}{probe}-derived truth directions and \glsterm{sparse-autoencoder}{SAE}-derived truth-related
features partially overlap, partially do not, and the cases of
disagreement are exactly the cases where a single-method claim should
not be taken as canonical.
